\vspace{-0.3cm}\section{Discussion and Outlook}\vspace{-0.3cm}
\label{sec:discussion}\begin{sloppypar}
Building LACK showed that valuable info can be automatically extracted, integrated and enriched with semantics.
Connections that are invisible within any single source become queryable once the sources are represented in a common graph, and this cross-source bridging is precisely what practitioners told us they value most.
We see no reason this approach should remain specific to climate lobbying.
The same approach, harvesting heterogeneous investigative documentation, extracting typed relations with large language models, and grounding entities against open knowledge bases, could in principle support accountability work in other domains, such as public-health lobbying or defence procurement.
Demonstrating this generalisation is an aim of our ongoing work.\end{sloppypar}

At the same time, the construction process surfaces challenges that are characteristic of this domain.
Lobbying networks are dynamic: positions, affiliations, and roles change over time, yet temporal information is the dimension our evaluation and our experts both found hardest to capture, and entity-level activity bounds remain sparse.
Funding is the relation practitioners most want and one of the hardest to extract reliably, since funding ties are often disclosed obliquely or not at all, as our experts also noted.
The long tail of minor actors, obscure consultants, short-lived coalitions, regional trade bodies, drives the high proportion of entities without external identifiers; rather than a mere limitation, this is an open research problem, since enriching and disambiguating long-tail entities absent from general-purpose knowledge bases is an active challenge for automatic knowledge graph construction~\cite{graciotti2025lost}.
The reliance on English-language sources introduces a geographic and linguistic skew, visible in our COP30 cross-reference, where actors from lower-income countries are markedly under-represented.

Turning a resource into something practitioners can adopt raises questions beyond extraction quality.
LACK reorganises information produced by investigative organisations who invest considerable effort in it, and respecting that effort is of primary importance: we followed DeSmog's republishing guidelines and complied with the redistribution terms we obtained from InfluenceMap, and we publish LACK under a Creative Commons Attribution-NonCommercial licence.
In an era when large commercial actors routinely ingest others' content without credit, a BY-NC licence preserves attribution and reserves commercial rights to the original data owners, which we hope makes it easier, not harder, for further investigative organisations to contribute to a shared LACK data space.
The gaps our experts identified, direct links to the political sphere, finer distinctions between intermediary actor types such as consultancies and PR agencies, and integration of further sources such as the EU Transparency Register and OpenSecrets, define our agenda for that growth.

In summary, we have presented LACK, to our knowledge the first openly published knowledge graph that turns fragmented investigative records on climate lobbying into a queryable resource, demonstrated its analytical value, and assessed both its quality and its reception among the practitioners it is meant to serve.
By releasing LACK as an open, attribution-preserving resource, we hope to invite the organisations whose work it builds on to help extend it, growing it from a knowledge graph into a shared data space for climate accountability.